% The glossary: every word of the book, defined once. A book's glossary lists only the words it
% uses (\term{key}, \termas{key}{text}), so each band's glossary is complete on its own.
% Library README: lib/README.md.

\newglossaryentry{active-pieces}{name={active pieces (strategy)},text={active pieces},
  description={Pieces that attack, control many squares and support their own pawns, instead of only defending. A passive piece often becomes a weakness itself.}}
\newglossaryentry{back-rank-mate}{name={back-rank mate},text={back-rank mate},
  description={Mate on the edge of the board: the rook checks along the last rank, and the king cannot leave it, because the enemy king or its own pieces take away the squares in front of it.}}
\newglossaryentry{barrier}{name={barrier (strategy)},text={barrier},
  description={A piece, often supported by pawns, controls a line of squares that the enemy king cannot cross.}}
\newglossaryentry{breakthrough}{name={breakthrough (king and pawn endings)},text={breakthrough},
  description={A pawn advance, usually with a sacrifice or an exchange, that opens the way for a pawn through the opponent's pawns and makes it a passed pawn.}}
\newglossaryentry{bridging}{name={bridging (rook endings)},text={bridging},
  description={The attacker puts his rook between the defending rook and his own king and pawn, so that the checks stop and the pawn can queen.}}
\newglossaryentry{centralize}{name={centralize the pieces (strategy)},text={centralize the pieces},
  description={Pieces in the centre control more squares and can go quickly to either wing. In queen endings, the queen in the centre can give or stop checks from everywhere.}}
\newglossaryentry{diagonal-checks}{name={check along the diagonal (queen endings)},text={check along the diagonal},
  description={A queen check along a diagonal. In queen endings, checks along the diagonals often lead to perpetual check more easily than checks along files and ranks.}}
\newglossaryentry{clear-cut-draw}{name={clear-cut draw (practical play)},text={clear-cut draw},
  description={A line that reaches a plainly drawn position by force, so that nothing is left to calculate or to go wrong.}}
\newglossaryentry{counter-check}{name={counter-check (queen endings)},text={counter-check},
  description={A check answered by a check that also blocks it, often forcing the exchange of queens.}}
\newglossaryentry{counterplay}{name={counterplay (practical play)},text={counterplay},
  description={Threats of your own, such as a passed pawn, that force the opponent to react instead of carrying out his plan.}}
\newglossaryentry{counting}{name={counting (practical play)},text={counting},
  description={Instead of playing a race out move by move, count how many moves each side needs, and compare. Tracing the moves one by one takes concentration and easily goes wrong by one tempo, and the same position with the other side to move can have the opposite result. Counting gets you to the final position quickly, and saves energy.}}
\newglossaryentry{cut-off}{name={cut off (rook endings)},text={cut off},
  description={A rook that controls a whole rank or file cuts off the enemy king: the king cannot cross that line. The defender uses it to keep the attacking king away from its pawn, the attacker to keep the defending king away from the pawn's path.}}
\newglossaryentry{dangerous-corner}{name={dangerous corner (rook vs bishop)},text={dangerous corner},
  description={In rook against bishop, the corner of the same colour as the bishop. There the defending king can be mated; the safe corner is the one of the other colour.}}
\newglossaryentry{desperado}{name={desperado},text={desperado},
  description={A piece that is lost anyway, or that has nothing better to do, gives itself up as dearly as possible: it captures or checks until it is taken. Often taking it is stalemate, or not taking it allows perpetual check.}}
\newglossaryentry{do-not-block-queen}{name={do not block your own queen (queen endings)},text={do not block your own queen},
  description={The king should not stand on the lines its own queen needs for checks or defence.}}
\newglossaryentry{do-not-lose-tempo}{name={do not lose a tempo (pawn races)},text={do not lose a tempo},
  description={In a pawn race every tempo counts: a single move that does not bring a pawn closer to its queening square, or does not stop the enemy pawn, can decide the race.}}
\newglossaryentry{do-not-rush}{name={do not rush (strategy)},text={do not rush},
  description={When the opponent has no real counterplay, there is no hurry. First take away his resources and improve your pieces, then act.}}
\newglossaryentry{do-not-weaken}{name={Do not weaken your own position (strategy)},text={Do not weaken your own position},
  description={A pawn move cannot be taken back. If the position holds as it is, leave the pawns where they are: every unnecessary pawn move can create a weakness or a passed pawn for the opponent.}}
\newglossaryentry{domination}{name={domination (strategy)},text={domination},
  description={A piece is attacked and has no safe square to go to.}}
\newglossaryentry{drive-off-the-bishop}{name={driving off the bishop (bishop endings)},text={driving off the bishop},
  description={Attacking the enemy bishop, or offering an exchange, so that it must leave the diagonal on which it does its job.}}
\newglossaryentry{easy-choice}{name={easy choice (practical play)},text={easy choice},
  description={A choice between a few moves, one of which wins, that the opponent can settle without much calculation. Avoid offering one. You can find the move that leaves no such choice without calculating which option would win, and so save your energy for the decisions that really need calculation.}}
\newglossaryentry{eliminate-the-pawn}{name={eliminate the pawn (rook endings)},text={eliminate the pawn},
  description={Win or exchange the attacker's last pawn. Without it, a rook ending is a draw at once.}}
\newglossaryentry{fortress}{name={fortress (strategy)},text={fortress},
  description={A defensive set-up that the stronger side cannot break, even with extra material. The defender only has to keep it.}}
\newglossaryentry{frontal-checks}{name={frontal checks (rook endings)},text={frontal checks},
  description={The defending rook checks the attacking king from the front, so that the pawn cannot move forward. Otherwise the defence is lost.}}
\newglossaryentry{interference}{name={interference (strategy)},text={interference},
  description={A piece is put on a line to cut it, so that an enemy piece can no longer guard a square through it.}}
\newglossaryentry{key-squares}{name={key squares (king and pawn endings)},text={key squares},
  description={The squares from which the king of the stronger side wins by force, whatever the defender does and whoever is to move.}}
\newglossaryentry{king-activity}{name={king activity (strategy)},text={king activity},
  description={In the endgame the king is a strong piece. Bring it to the centre, or to where the play is, before the opponent's king gets there.}}
\newglossaryentry{king-after-king}{name={king after king (bishop endings)},text={king after king},
  description={The defending king follows the attacking king and stands opposite it, so that the attacker's bishop cannot interfere with the defending bishop's diagonal.}}
\newglossaryentry{long-side}{name={long side (rook endings)},text={long side},
  description={The wider of the two parts into which a pawn (other than a rook's pawn) divides the board. The defending rook belongs there, so that it has room for side checks.}}
\newglossaryentry{lucena}{name={Lucena's position (rook endings)},text={Lucena's position},
  description={The attacker's pawn is on the seventh rank with its king in front of it, and the defending king is cut off. The attacker wins by bringing his rook to the fourth rank, walking his king out, and blocking the checks from behind with the rook.}}
\newglossaryentry{mutual-zugzwang}{name={mutual zugzwang (king and pawn endings)},text={mutual zugzwang},
  description={A position in which the side to move, whichever it is, would rather pass.}}
\newglossaryentry{one-diagonal}{name={one-diagonal principle (bishop endings)},text={one-diagonal principle},
  description={A bishop defends well when everything it has to do lies on one diagonal. When it must watch two diagonals at once, it can be torn apart.}}
\newglossaryentry{opposite-corner}{name={opposite corner (queen endings)},text={opposite corner},
  description={Against a knight's or rook's pawn, the corner diagonally opposite the pawn's queening square, also called the drawing zone. There the defending king is safest, because the attacker can hardly answer a check with a counter-check.}}
\newglossaryentry{opposition}{name={opposition (king and pawn endings)},text={opposition},
  description={The kings stand on the same file (or rank) with one square between them, and the other side is to move, so its king has to give way.}}
\newglossaryentry{pawns-other-colour}{name={other colour than the bishop (bishop endings)},text={other colour than the bishop},
  description={Put your pawns on squares of the colour your bishop does not move on. Then bishop and pawns together control squares of both colours, the bishop does not have to guard its own pawns, and the pawns do not block its diagonals.}}
\newglossaryentry{outflanking}{name={outflanking (king and pawn endings)},text={outflanking},
  description={Going round the side of the enemy king when it has given way, to reach the key squares.}}
\newglossaryentry{passive-defence}{name={passive defence (rook endings)},text={passive defence},
  description={The defending king stays in front of the pawn and the rook waits on the first rank. It holds against a knight's pawn, but loses against a bishop's pawn or a central pawn.}}
\newglossaryentry{perpetual-check}{name={perpetual check (queen endings)},text={perpetual check},
  description={A series of checks the other king cannot escape. The game is drawn.}}
\newglossaryentry{philidor}{name={Philidor position (rook endings)},text={Philidor position},
  description={The defending king stands in front of the pawn, and the defending rook on its third rank keeps the attacking king away. To make progress the attacker has to advance the pawn, and then the rook goes to the far end of the board and checks from behind.}}
\newglossaryentry{ponziani}{name={Ponziani position (bishop endings)},text={Ponziani position},
  description={A king holds against bishop and knight's pawn: the king stands in front of the pawn, and the bishop is shut in the corner beside it, for example a black bishop on h2 and pawn on g3 against a white king on g2. The bishop can never get out, and the game is a draw, wherever the other king stands.}}
\newglossaryentry{process-of-elimination}{name={process of elimination (practical play)},text={process of elimination},
  description={Ruling out a candidate move because the opponent has a simple reply to it: when you are playing for a win, a simple reply that draws; when you are playing for a draw, a simple reply that wins for him. The move needs no further calculation.}}
\newglossaryentry{prophylaxis}{name={prophylaxis (strategy)},text={prophylaxis},
  description={Preventing the opponent's plan before it can start. Dvoretsky made it a habit of thought: ask what the opponent would play if it were his move, and take that possibility away, for example his only counterplay, instead of only pursuing your own plan.}}
\newglossaryentry{rear-checks}{name={rear checks (rook endings)},text={rear checks},
  description={The defending rook checks the attacking king from behind, from the far end of the board, along the files. The king cannot hide in front of its pawn, because the defending king stands there.}}
\newglossaryentry{remote-passed-pawn}{name={remote passed pawn (king and pawn endings)},text={remote passed pawn},
  description={A passed pawn far away from the other pawns. It lures the enemy king away, while one's own king attacks the pawns on the other side.}}
\newglossaryentry{removing-the-pawns}{name={removing the pawns (strategy)},text={removing the pawns},
  description={The defender wins or exchanges the stronger side's pawns, giving up a piece if necessary. Without pawns, the extra piece cannot win.}}
\newglossaryentry{rook-against-pawn}{name={rook against pawn (rook endings)},text={rook against pawn},
  description={An ending in which one side has only a rook and the other only a pawn, usually after a rook was given up for a pawn. The rook wins when its king is close enough or the defending king is cut off from the pawn; the pawn has drawing chances only when its own king supports it and the other king is far away.}}
\newglossaryentry{rule-of-the-square}{name={rule of the square (king and pawn endings)},text={rule of the square},
  description={If the defending king can step into the square from the pawn to its queening square, it catches the pawn; if it cannot, the pawn queens.}}
\newglossaryentry{schematic-thinking}{name={schematic thinking (practical play)},text={schematic thinking},
  description={Planning by picturing the set-up you want and the piece moves that lead there, instead of calculating every reply of the opponent.}}
\newglossaryentry{short-side}{name={short side (rook endings)},text={short side},
  description={The narrower of the two parts into which a pawn (other than a rook's pawn) divides the board. The defending king belongs there, where it does not stand in the way of its own rook.}}
\newglossaryentry{side-checks}{name={side checks (rook endings)},text={side checks},
  description={The defending rook checks the enemy king from the side, along the ranks. The defence holds when the king cannot escape the checks without leaving its pawn.}}
\newglossaryentry{simple-logic}{name={simple logic (practical play)},text={simple logic},
  description={Reasoning from results you already know instead of calculating moves: if every way to make progress leads to a position you know is a draw, the position is a draw.}}
\newglossaryentry{skewer}{name={skewer},text={skewer},
  description={An attack along a line on a piece that must move away, uncovering the piece behind it, which is then taken.}}
\newglossaryentry{static-advantage}{name={static advantage (strategy)},text={static advantage},
  description={A lasting advantage, such as a remote passed pawn or a better pawn structure, that does not go away by itself. The other side must react at once, dynamically, and create counterplay; otherwise it loses in the long run.}}
\newglossaryentry{bishop-escorts}{name={the bishop escorts the pawn (bishop endings)},text={the bishop escorts the pawn},
  description={The bishop supports the pawn's advance by guarding the squares in front of it, so that the king is free for other tasks.}}
\newglossaryentry{king-escorts}{name={the king escorts the pawn (king and pawn endings)},text={the king escorts the pawn},
  description={The king walks next to or in front of its pawn and guards the squares on its way, so that the pawn can queen even with the enemy king close by.}}
\newglossaryentry{three-file-rule}{name={three-file rule},text={three-file rule},
  description={A rook can only give lasting checks from a distance: at least three files (or ranks) away from the attacking king and pawn. Closer than that, the king walks towards the rook and the checks run out.}}
\newglossaryentry{tie-down}{name={tie down (strategy)},text={tie down},
  description={Attack a weakness, or pin a piece, so that an enemy piece must stay where it is to defend. A tied-down piece cannot take part in the play, and when it leaves, the weakness falls.}}
\newglossaryentry{transposing}{name={transposing into a pawn ending (strategy)},text={transposing into a pawn ending},
  description={Exchanging the last pieces to reach a pawn ending whose result you can judge. Pawn endings can often be judged exactly, by counting or by rules such as Bähr's.}}
\newglossaryentry{undermining}{name={undermining (strategy)},text={undermining},
  description={Attacking the pawn that supports another pawn, so that the supported pawn loses its protection and becomes weak.}}
\newglossaryentry{unnecessary-check}{name={unnecessary check (practical play)},text={unnecessary check},
  description={A check that achieves nothing. It often drives the enemy king to a better square, or costs a tempo.}}
\newglossaryentry{wrong-bishop}{name={wrong bishop (bishop endings)},text={wrong bishop},
  description={A bishop that does not control the promotion square of a rook's pawn. With such a bishop and a rook's pawn, the stronger side cannot drive the defending king out of the corner, and the game is a draw.}}
\newglossaryentry{zugzwang}{name={zugzwang},text={zugzwang},
  description={The side to move would rather pass: every move makes its position worse.}}
